\section{Background and Related Work}
\label{sec:related}

\subsection{High-Latency Transmission and End-to-End Control}
High-latency transmission has long been studied in satellite and other high-bandwidth-delay-product networks. Standard TCP mechanisms can be extended with window scaling, selective acknowledgements, timestamps, and congestion-control improvements, but the sender still depends on delayed end-to-end feedback. Model-based and delay-aware algorithms such as BBR and Copa improve the trade-off between throughput and delay by estimating path properties or optimizing utility objectives. Learning-based congestion control further explores online adaptation under changing environments. These studies provide important reference points, but their control decisions are still made at the endpoint. In dynamic high-RTT links, the feedback delay can remain a fundamental obstacle when bandwidth and loss conditions switch within a time scale comparable to or shorter than the RTT.

\subsection{Performance Enhancing Proxies}
PEPs were proposed to mitigate link-related performance degradations by introducing intermediate processing such as connection splitting, local acknowledgement, caching, compression, and protocol conversion. In high-latency networks, PEPs are valuable because they move part of the control and recovery process closer to the bottleneck link. PEPesc is the direct baseline of this work. It deploys distributed proxy entities around the high-latency link and combines proxy-side bandwidth estimation with streaming coding recovery. This design reduces the dependence on RTT-scale retransmission and provides a practical basis for proxy-enhanced high-latency transmission. However, when bandwidth and loss rate vary rapidly, the original rule-based perception and control logic may not react quickly enough, leaving room for dynamic perception and control enhancement.

\subsection{Explicit Feedback and Local Control}
Network-side explicit feedback approaches, including XCP, RCP, and ABC, shorten the control loop by allowing bottleneck-side or network-side elements to participate in rate adjustment. ABC is particularly relevant because it abstracts control behavior into accelerate and brake semantics under time-varying wireless links. Such methods inspire the design of interpretable control states. Nevertheless, they typically require support from routers or network devices. In contrast, this paper assumes that only the proxy nodes can be controlled. Therefore, we construct a proxy-internal state machine that mimics the clarity of explicit feedback while relying only on locally observable proxy-side signals.

\subsection{Streaming FEC and Loss Recovery}
In high-RTT networks, loss recovery through retransmission incurs at least one additional feedback loop and can significantly delay application-layer delivery. FEC and on-the-fly erasure coding reduce this dependence by injecting redundancy before losses are fully known. PEPesc adopts streaming FEC to recover losses continuously during transmission. This mechanism is especially important when random loss is present. However, redundancy control must be coordinated with pacing: too little redundancy may interrupt decoding and feedback, whereas excessive redundancy consumes bottleneck bandwidth and can worsen queue buildup. This trade-off motivates the FEC reliability floor in our execution layer.

\subsection{Learning-Assisted Network Perception}
Network states such as bandwidth, RTT, ACK timing, queue delay, and loss evolve over time and exhibit temporal dependency. Recurrent neural networks, especially GRU and LSTM variants, are suitable for modeling such time series. In system and networking tasks, however, learning components must be deployed conservatively because inaccurate predictions can destabilize the control loop. This paper therefore uses GRU only at the perception layer. The model learns residual corrections over a rule-based estimator, while the final control decisions remain constrained by explicit states, conservative pacing correction, and FEC protection.
